%======================================================================
\section{Testing every model on real data}
\label{sec:realtest}

\begin{basic}
The synthetic world only ranks methods. On the real CCD the models are tested in stages, each stage using data that the packets
produce anyway, and each with a test that a model cannot pass by memorizing. Stage 0 needs no training at all: it checks whether the
physics that model A relies on holds on our CCD.
\end{basic}

\paragraph{Stage 0 -- the first packets at nominal settings (no training).}
\begin{itemize}
\item \textbf{Label noise.} Repeat the nominal packet $\gtrsim$5 times: the scatter of WADERs' $\sigma_e$, $\lambda$ and gain between repeats is the
  floor no model can beat (synthetic: 1.5\,\%, 3\,\%, 0.8\,\%). Compare with WADERs' own error bars (\code{sigma\_e\_err}, \code{lambda\_err}).
\item \textbf{Does the trace physics hold?} With the gain from the image, the trace alone predicts
  $\sigma_e^\text{trace} = \sqrt{a^2/N + b^2}/g$ ($a, b$ from the Allan deviation, or the features ``skip: CDS rms'' and
  ``fit: correlated noise floor''). The ratio $\sigma_e^\text{image}/\sigma_e^\text{trace}$ and its scatter decide how much model A must
  learn: a ratio near 1 with small scatter means A works almost without training (world v2); a ratio that depends on the settings
  is what A's residual $r(x)$ learns (world v3).
\item \textbf{Dark current vs CIC.} Fit $\lambda$ row by row (16 rows): the slope is the dark current accumulated while the rows wait, the
  intercept the CIC. A clear slope means the per-row labels are worth adding to the dataset.
\end{itemize}

\paragraph{Stage 1 -- the first scan.} All models of both rounds, the same splits and scores:
\begin{itemize}
\item \textbf{splits:} (i) grouped by packet (5 folds); (ii) \emph{blocked in time}: train on the first 2/3 of the scan, test on the last
  1/3 (does the model follow the drift?); (iii) \emph{leave one parameter out} for a one-at-a-time scan: train on 9 parameters' sweeps
  and the references, test on the 10th (it measures what the model can \emph{not} know, and is the honest test for a design without joint
  changes);
\item \textbf{scores:} relative error against the labels, compared with the label noise of stage 0; coverage of the 68\,\% and 95\,\%
  intervals before and after conformal calibration (H); for the surrogate, the rank correlation between the predicted and measured
  objective and the regret of the best predicted point;
\item \textbf{learning curves:} the same, on random subsets of 15, 30, 60, 120, 240 packets;
\item \textbf{what the models use:} GP relevances, A's residual $r(x)$ (where the trace physics fails), and, for the CNN, which parts of the
  curves matter (saliency).
\end{itemize}

\paragraph{Stage 2 -- a second module (when one is available).} Train on module 1 plus $k = 0, 1, 2, 5, 10, 20$ packets of module 2;
test on its other packets (the \code{newmodule} experiment). This decides B against the first round and how many packets a new module needs.

\paragraph{Stage 3 -- the optimizer.} First \emph{offline}: replay Bayesian optimization on the scan, choosing only among the measured
packets (pool-based), for every surrogate; cheap and fair. Then one \emph{prospective} run with the best surrogate: it proposes, we take the
packet, and the prediction made before the packet is compared with the result. Only this last test measures what we care about.

\paragraph{Stage 4 -- the data-hungry models (C, D, E, F).} Once there are tens of thousands of traces: pre-train the encoder on all of them
(D), put it under the GP (C), and compare with A+B at the same number of labeled packets. Set models (E) need the individual trace sets
(see the data format below).

%======================================================================
\section{Do the packets need to change?}
\label{sec:dataneeds}

\begin{basic}
Mostly no: the packet as designed (5 reference sets, 10 trace sets, one 16-row image) carries what every model needs. Four additions
would make the models clearly better, and one choice in the scan design matters a lot. None of them makes a packet longer.
\end{basic}

\begin{table}[h]
\centering\small
\begin{tabularx}{\textwidth}{>{\raggedright\arraybackslash}p{4.3cm}>{\raggedright\arraybackslash}X>{\raggedright\arraybackslash}p{2.4cm}}
\toprule
Change & Why (which model) & Cost \\
\midrule
\textbf{Per-row $\lambda$ from WADERs} (and per-row $\sigma$) & splits dark current from CIC (A); 16 constraints per image instead of 1; was already an
  open item (WADERs JSON with per-row dark current) & none (analysis only) \\
\textbf{Keep the un-averaged skips of the image}, or of a few of its rows, if the DAQ can save them & the image's own $\sigma(N)$ curve: $\sigma_1$,
  floor and the image-only noise measured on the image, not inferred (A; stage 0) & disk (tens of MB per channel for 16 rows; check that cdaq can save them) \\
\textbf{A scan design with joint changes}: random/Sobol points inside the limits, or half one-at-a-time and half random & the
  one-at-a-time design never changes two parameters together, so no model can learn interactions (the \code{design} experiment, below) & none
  (same number of packets) \\
\textbf{Trace-only packets} between the image packets, at extra settings & unlabeled traces for pre-training (D), the cheap fidelity of
  multi-fidelity BO (F), and many more $\theta$ points for the trace physics $a(\theta), b(\theta)$ (A) & minutes each, no image \\
\textbf{Record the conditions}: CCD and board temperatures, time since the last flush/erase, vacuum, the complete parameter set of every pin
  (not only the scanned ones), software versions & the drifting state is today inferred from the reference traces; measured temperatures
  make it direct (all models; B for modules) & none \\
\textbf{Nominal image packets at intervals} (the scan already has \code{reference\_every = 5}) & labeled drift anchors: they measure the
  state's effect on the labels directly & already planned \\
\textbf{A few packets of every new module}, at nominal and around it & B needs $\gtrsim$2--5 imaged packets per new module before its gain
  (and so $\sigma_e$) is known; traces alone cannot give the gain & 2--5 packets \\
\bottomrule
\end{tabularx}
\caption{Changes to the data that the second-round models would use. None makes a packet longer.}
\label{tab:dataneeds}
\end{table}

%======================================================================
\section{A better data format}
\label{sec:format}

\begin{basic}
Today the models read \code{dataset.csv}: one row per packet and amplifier, every feature already averaged over the repeats. That is
the right table for the first-round models, but it throws away what the next ones need: the individual repeats, the curves, the label
error bars, the row-by-row image information and the conditions. The proposal is to keep the raw files as they are and add a small set
of linked tables, from which \code{dataset.csv} is just one view.
\end{basic}

\begin{advanced}
\textbf{Keep} the raw data as the source of truth (FITS images, lossless \code{.npz} traces, \code{result.json}). \textbf{Add} a tidy,
columnar store (Parquet tables plus one array store, HDF5 or Zarr), keyed by \code{packet\_id}, \code{board}, \code{channel}, \code{kind}
(reference / $\theta$) and \code{set}:
\begin{enumerate}
\item \code{packets}: time stamps, module and CCD identifiers, configuration hashes, the complete parameter vector (every pin),
  temperatures and other conditions, readout times, flags; also the trace-only packets;
\item \code{trace\_sets}: one row per trace set (not combined): quality flags and all 85 features, so that set models (E) and robust
  combinations can be learned, and the repeat scatter is a feature in its own right;
\item \code{curves}: the 12 curves of every trace set, at their natural length (array store), for the CNN, pre-training (D) and deep kernels (C);
\item \code{labels}: per packet and amplifier $\sigma_e$, $\lambda$, gain \emph{with their errors} (for heteroscedastic likelihoods: the GP
  then uses WADERs' error instead of learning one noise level), fit flags, WADERs version;
\item \code{image\_rows}: per row (and column) $\lambda$, $\sigma$, baseline: the information model A uses; and, if available, the image's own
  $\sigma(N)$ curve;
\item \code{schema}: version and the hash of the code that produced every table.
\end{enumerate}
\code{dataset.csv} becomes a derived view (combine the sets, join the labels), so the first-round code is unchanged. Trace-only packets
are rows without labels; new modules are new identifiers. Loading one table into pandas or PyTorch is a single call.
\end{advanced}
